\chapter{1993 Nobel Prize in Economics}
\textbf{Laureates: }Robert Fogel (USA), Douglass North (USA)

\section*{Reason for Award}
New economic history quantitative analysis.

\section{What They Did}
Robert Fogel pioneered cliometrics applying quantitative methods to historical economic analysis challenging established narratives about economic growth Slavery Profitability demonstrating slave plantation profitability contrary prevailing assumptions Railway contribution assessment showing limited productivity gains Alternative interpretations prompted rigorous reexamination of historical economic facts Douglass North developed institutional economics analyzing how institutions shape economic performance over time Created formal framework distinguishing formal rules informal constraints implementation characteristics Analyzed path dependency in institutional development Examinined role ideology beliefs in shaping institutional trajectories Historical institutionalism examined how small initial events large long-run consequences through lock-in effects Together Fogel North demonstrated importance careful empirical historical examination combined theoretical rigor to resolve longstanding economic historical debates

\section{Impact on Economic Thinking}
Fogel and North transformed economic history from descriptive narrative discipline scientific inquiry grounded in theory and evidence Cliometrics established demanding standards requiring explicit counterfactuals rigorous measurement statistical inference North s institutional framework provided explanatory power for divergent national development paths Their collaborative emphasis combining quantitative rigor theoretical coherence raised expectations for historical research Institutions developed as critical variable explaining cross-country income differences Path dependency concepts revealed persistence of institutions over centuries Fogel s provocative findings challenged conventional wisdom promoting skeptical scrutiny of historical claims Historians adopted cliometric methods producing richer more nuanced understanding past Economists incorporated historical perspective avoiding ahistorical abstraction Their joint contributions created new interdisciplinary field combining economic theory historical evidence quantitative methods institutional analysis.

\section{Modern Perspectives Today}
Cliometrics remains essential economic historical research methodology National account reconstruction historical estimates employ Fogel-style counterfactual analysis Institutional economists study persistence change following North s frameworks Comparative historical analysis examines institution origins diffusion effectiveness Historical development studies apply path dependency concepts Understanding colonial legacies governance transitions ethnic conflicts all rely on institutional analytical frameworks Economic historians test institutional hypotheses using historical natural experiments Big digitization of historical records enables richer cliometric studies New institutional development programs incorporate North s insights about inclusive versus extractive institutions Historical institutionalism explains contemporary policy resistance to reform despite clear inefficiencies Fogel North legacy visible every historical economic analysis examining root causes prosperity poverty their institutional quantitative synthesis provides essential tools understanding economic evolution over time
